\documentclass[11pt,a4paper]{article}
\usepackage[T1]{fontenc}
\usepackage[utf8]{inputenc}
\usepackage[english]{babel}
\usepackage{amsmath,amssymb}
\usepackage[margin=2cm]{geometry}
\setlength{\parindent}{0pt}
\setlength{\emergencystretch}{3em}
\allowdisplaybreaks

\title{False identities with a declared metric (18)}
\author{}
\date{}
\begin{document}
\maketitle
Each statement is a true identity of \texttt{metric\_identities.tex} with one small change: a sign, a factor, the signature of the metric, the declared value of $\varepsilon$, or two exchanged indices. None of them holds.

\begin{enumerate}
% The sign of the determinant, the signature, the dimension.
\item $\displaystyle g_{ij} = \operatorname{diag}(-1, 1, 1, 1) \land \varepsilon_{1234} = 1 \implies \varepsilon^{ijkl}\varepsilon_{ijkl} = 24$
\item $\displaystyle g_{ij} = \operatorname{diag}(1, 1, 1, 1) \land \varepsilon_{1234} = 1 \implies \varepsilon^{ijkl}\varepsilon_{ijkl} = -24$
\item $\displaystyle g_{ij} = \operatorname{diag}(1, -1, -1, -1) \land \varepsilon^{1234} = 1 \implies \varepsilon_{1234} = 1$
\item $\displaystyle g_{ij} = \operatorname{diag}(-1, 1, 1, 1) \land \varepsilon_{1234} = 1 \implies \varepsilon^{ijkl}\varepsilon_{ijkm} = 6\delta^{l}_{m}$
\item $\displaystyle g_{ij} = \operatorname{diag}(-1, 1, 1, 1) \land \varepsilon_{1234} = -1 \implies \varepsilon^{1234} = -1$
\item $\displaystyle g_{ij} = \operatorname{diag}(-1, 1) \land \varepsilon_{12} = 1 \implies \varepsilon^{ij}\varepsilon_{ij} = 2$
\item $\displaystyle g_{ij} = \operatorname{diag}(1, 1, 1) \land \varepsilon_{123} = 1 \implies \varepsilon^{ijk}\varepsilon_{ijk} = -6$
% The metric and its inverse.
\item $\displaystyle g_{ij} = \operatorname{diag}(-1, 1, 1, 1) \implies g^{ij} g_{ij} = 2$
\item $\displaystyle g_{ij} = \operatorname{diag}(-1, 1, 1, 1) \implies \forall i, k \in \{1, 2, 3, 4\}\ \left( g^{ij} g_{jk} = 2\delta^{i}_{k} \right)$
\item $\displaystyle g_{ij} = \operatorname{diag}(2, 3) \implies g^{ij} g_{ij} = 5$
\item $\displaystyle g_{ij} = \operatorname{diag}(2, 3) \implies g_{11} + g_{22} = 2$
\item $\displaystyle g^{ij} = \operatorname{diag}(2, 3) \implies g_{11} + g_{22} = 5$
% Raising and lowering.
\item $\displaystyle g_{ij} = \operatorname{diag}(-1, 1, 1, 1) \implies \forall i \in \{1, 2, 3, 4\}\ \left( g^{ij} A_{j} = 2 A^{i} \right)$
\item $\displaystyle g_{ij} = \operatorname{diag}(-1, 1, 1, 1) \implies A^{i} B_{i} = -A_{i} B^{i}$
\item $\displaystyle g_{ij} = \operatorname{diag}(-1, 1, 1, 1) \implies \forall i, j \in \{1, 2, 3, 4\}\ \left( F^{ij} = g^{ik} g^{jl} F_{lk} \right)$
\item $\displaystyle g_{ij} = \operatorname{diag}(-1, 1, 1, 1) \implies g_{ij} A^{i} B^{j} = g_{ij} A^{i} A^{j}$
\item $\displaystyle g_{ij} = \operatorname{diag}(-1, 1, 1, 1) \implies T^{i}{}_{i} = 2 g^{ij} T_{ij}$
\item $\displaystyle g_{ij} = \operatorname{diag}(2, 3) \implies g_{ij} A^{i} A^{j} = 2 A^{i} A_{i}$
\end{enumerate}
\end{document}
